\section{Introduction}
\label{sec:intro}
Strong game-playing systems spend roughly the same computation on every move, although some decisions are plainly harder than others. Adaptive allocation of search has a long history, from selective extensions in alpha--beta engines~\cite{knuth1975,campbell2002} to time management for Monte Carlo tree search (MCTS)~\cite{huang2010} and the metareasoning literature~\cite{russell1991,tolpin2012}. Dual-process agents separate a fast learned ``intuition'' (System~1) from slower deliberative search (System~2)~\cite{anthony2017,booch2021,kahneman2011,evans2013}, which raises a \emph{gating} question: when should System~2 be invoked?

A tempting answer is to let System~1 decide: if the fast model is unsure, search. This assumes predictive uncertainty is a good proxy for the \emph{value of additional computation}. The two need not coincide. A position can be easy to rank yet need search to avoid a refutation, and a position can look unclear while all plausible moves are equally good, so search has nothing to find. We test the assumption directly.

\paragraph*{What we built and measured}
System~1 is LayaChess~\cite{layachess}, a 421M-parameter model that returns a win chance for every legal move. System~2 is a PUCT search~\cite{kocsis2006,rosin2011,silver2017} seeded by System~1's scores, with leaves evaluated by a 2.3M-parameter value network we trained. A gate converts System~1's scores into a confidence and searches when it is low. Stockfish~\cite{stockfish} serves only as judge of move quality and rating anchor. The design follows a protocol committed before the test runs (Section~\ref{sec:hyp}), and its central device is a \emph{matched random gate}: a control that searches the same fraction of positions by coin flip. A gate that cannot beat its random control has not shown that confidence carries information about \emph{when} search is useful.

\paragraph*{How the study progressed}
Stage V1 answered the gating question negatively. We then continued in two stages on fresh test months so that earlier tuning could not leak into the test. V2 asked whether repairing the weaknesses V1 exposed makes the agent stronger; V3 asked again whether a gate can beat a random gate, using a learned gate, and whether System~1 contributes once the repairs exist (Section~\ref{sec:followup}). V1 is preserved as the confirmatory study; the later stages test different agents under different conditions and are not pooled with it.

\paragraph*{Contributions}
(1) A protocol-driven evaluation of confidence-based selective search with matched random-gate controls: \nPos{} judged positions, a 500-game anchored tournament and 1{,}000 puzzles. (2) A negative primary result: the confidence gate is indistinguishable from a random gate, and two further predictions are reversed; a post hoc analysis suggests confidence tracks how dominant one move is, not how much search would help. (3) A follow-up showing that a shallow material blunder filter lowers ACPL by 38\cp\ and lifts Elo by about 200, while a learned gate beats random gating only modestly (about 2.5\cp), significantly only under a post hoc test. (4) A statistical audit: every V1 number is reproduced from raw data, as are all V2 and V3 statistics reported here, and a configuration error that invalidated a first V3 run is disclosed.
